\section{Embedding quantitative and analytic mathematics}

\subsection{Rational coding and error predicates}

The native domain is \(\N\), so signed integers, rationals, finite vectors, and
matrices are represented by primitive-recursive codes.  A rational may be a
pair \((z,d)\) with a signed-integer code \(z\) and positive denominator
\(d+1\).  Equality can be checked by cross multiplication, and order by the
corresponding signed comparison.  Approximation statements use integer-scaled
inequalities, avoiding division in the proposition itself.

For example, the assertion that rational code \(r\) approximates a real-like
object \(x\) within \(2^{-n}\) is represented by a certificate predicate
\[
 \mathsf{Approx}(x,n,r,c)=1,
\]
where \(c\) contains the finite evidence appropriate to the chosen
representation.  The framework never quantifies over a completed real number
as a set of rationals.  It reasons about finite observations at requested
precisions.

\subsection{Finite-use real interfaces}

A real-like object is supplied through a query contract
\[
 \forall n<N\;\exists c<B(x,n)\;
   \mathsf{Approx}(x,n,c)=1.
\]
For a theorem with an arbitrary precision parameter \(N\), the theorem is
proved uniformly in \(N\), but each instance uses finitely many queries.  A
reverse demand pass computes how accurate each input observation must be to
obtain the requested output precision.

This resolves a common apparent conflict: iPRA can prove a formula for every
natural precision parameter without constructing an infinitary object.  What
it cannot do automatically is derive the modulus that tells the finite
algorithm which input precision suffices.  The modulus is a primitive-
recursive producer and must be proved from theorem-specific estimates.

\subsection{Reverse error-demand propagation}

Consider a fixed finite computation graph.  Each node \(v\) has an output error
budget \(e_v\), and each edge carries a monotone local demand transformer
\(D_{u\leftarrow v}(e_v)\).  Starting from the requested final budget, propagate
requirements backward and take coordinatewise maxima.  If dependencies are
acyclic, one reverse traversal suffices.  If they are cyclic but bounded and
monotone, use the finite closure theorem of Section~7.

The forward proof then checks a local estimate at every node.  Typical local
contracts are
\begin{align*}
 |x+y-(\tilde x+\tilde y)|&\le e_x+e_y,\\
 |xy-\tilde x\tilde y|&\le M_xe_y+M_ye_x+e_xe_y,\\
 |f(x)-f(\tilde x)|&\le\omega_f(e_x),\\
 d(O(c),O(c'))&\le\varepsilon.
\end{align*}
All quantities are rational codes, and inequalities are reduced to natural
arithmetic.  Magnitude bounds \(M_x,M_y\) and moduli \(\omega_f\) are explicit
inputs or derived local lemmas.

\subsection{Division and apartness}

A bare negated equality is usually not a usable numerical resource for
approximate inversion.  Knowing \(x\ne0\) does not provide a rational
\(\lambda>0\) with \(|x|\ge\lambda\).  Without such a lower bound, an error in
\(x\) can be amplified arbitrarily by inversion.

\begin{contract}[Approximate inverse]
To compute an inverse with output error \(\varepsilon\), require:
\begin{enumerate}
\item an approximation \(\tilde x\) to \(x\);
\item a positive rational lower bound \(\lambda\le |x|\);
\item input error small enough, for example \(|x-\tilde x|\le\lambda/2\).
\end{enumerate}
Then \(|\tilde x|\ge\lambda/2\), and elementary rational algebra gives an
explicit inverse error bound.
\end{contract}

This requirement is not a weakness of iPRA; it is the constructive content of
numerical inversion.  In a discrete bounded natural-number problem the
situation is different.  If \(d<M\) and \(d\ne0\), finite numerical
case analysis yields \(1\le d\), which is a usable lower bound.  STRIVE keeps
these cases separate.

\subsection{Finite nets instead of abstract compactness}

An abstract compactness theorem is not directly a primitive-recursive
algorithm.  The useful finite contract is a net certificate:
\[
 \mathsf{Cover}(X,n,L)=1,
\]
where \(L\) is a finite list of sample points and every admissible input is
within radius \(2^{-n}\) of some member.  If a function has a supplied modulus
of continuity, evaluate it on the net, choose a least finite optimum, and add
the modulus error.

\begin{theorem}[Finite-net optimization schema]
Assume a checked \(\delta\)-net \(L\) for a domain, a bounded evaluator for
\(f\) on members of \(L\), and a modulus proving
\(d(x,y)\le\delta\to |f(x)-f(y)|\le\varepsilon\).  Then finite search returns a
sample whose value is within the evaluation error plus \(\varepsilon\) of the
global optimum.
\end{theorem}

\begin{proof}
Let \(x^*\) be an arbitrary admissible point.  Coverage supplies a sample
\(y\in L\) within \(\delta\).  The modulus compares \(f(x^*)\) and \(f(y)\).
The selected finite optimum is no worse than \(y\), up to its verified
evaluation error.  Combining the inequalities yields the stated bound.  The
sample selection and its failure alternative are instances of the signed
bounded-formula compiler.  If the theorem does not already provide a net or
modulus, those remain residuals.
\end{proof}

\subsection{Metastability and finite windows}

A full convergence modulus may not exist or may be hard to obtain.  Many finite
uses need only a metastable window:
\[
 \forall g<G\;\exists n<N(g,k)\;
 \forall i,j\in[n,n+g(n)]\;d(x_i,x_j)\le2^{-k}.
\]
When \(g\), \(k\), and a rank bound are finite-coded, this is a bounded formula
and receives positive or negative strategy certificates.  A positive strategy
returns a stable window for each bounded challenge.  A negative strategy
identifies a challenge and closes every candidate window with a violating pair.
This is a genuine constructive result even when no tail convergence rate is
claimed.

\subsection{Algebraic branches and clusters}

Exact branch selection is replaced by a hierarchy of finite contracts.

\begin{enumerate}
\item A local root certificate verifies that a finite box contains the claimed
number of roots or that a sign/argument condition holds.
\item A cluster observer returns symmetric data, such as coefficient
approximations or power sums, invariant under permutation.
\item A separation certificate upgrades the cluster to uniquely labeled
branches.
\item A compatibility trace is used only when a downstream theorem needs the
same branch across several precisions.
\end{enumerate}

The framework can verify and compose all these certificates.  It does not
produce a global separation theorem from the mere existence of roots.

\subsection{Diophantine enumeration}

Given a polynomial equation and a height cutoff \(H\), enumeration of all
integer tuples of height below \(H\) is a fixed bounded search.  STRIVE produces
both a solution certificate and a counterstrategy showing that every bounded
candidate fails.  What it cannot infer is a cutoff containing all solutions
from a non-effective finiteness assertion.

Thus the generic theorem has the conditional form
\[
 \mathsf{AllSolutionsBelow}(H)
 \;\longrightarrow\;
 \mathsf{EnumerationComplete}(H,L).
\]
The verifier and every consequence of a supplied \(H\) can be fully formalized.
The height producer is an exact, theorem-specific residual.
